%\documentclass{article}
%\usepackage{graphicx} % Required for inserting images

\title{Évaluation}
\author{Registia d'ALMEIDA}
\date{April 2026}

%\begin{minipage}{10cm}
%\vspace {-1cm}
%\begin{center}
%\textbf{\large RÉPUBLIQUE DU BENIN}\\
%\hspace{-0.5cm}\textbf{\large MINISTERE DE L'ENSEIGNEMENT SUPÉRIEUR}
%\textbf{\large SUPÉRIEUR ET DE LA RECHERCHE }
%\textbf{\large SCIENTIFIQUE(MESRS)\\
%\end{center}
%\end{minipage}




\begin{center}
\vspace{1cm}
    \textbf{I-\hspace{1cm}{MATHÉMATIQUES \hspace{0.3cm}$(7pts)$}}
\end{center}

\begin{center}
Soit $f$ une fonction continue périodique.On considère ses coefficients:
$$C_{n(f)}=\frac{1}{2\pi}\int_{2\pi}^{0} f(t)e^{-int},dt $$
\begin{enumerate}
    \item Prouver que 
    $\sum_{+\infty}^{-\infty}|C_{n}(f)\vert^{2}=\frac{1}{2\pi}\int_{0}^{2\pi}|f\rvert^{2},dt$
\end{enumerate}
\vspace{1cm}
\textbf{II-\hspace{1cm}{PHYSIQUE \hspace{0.3cm}$(7pts)$}}
\begin{center}


\end{center}
Démontrer l'invariant de la pseudo-norme du quadrivecteur impulsion :
$$(\frac{E}{c})^{2}-p^{2}=m^{2}c^{2}$$
\vspace{1cm}
Calculer l'enthalpie libre de réaction $\Delta_{r}G$ en fonction du quotient de réaction $Q$:
$$\Delta_{r}G=\Delta_{r}G^{O}+RT\ln(Q)$$

\begin{minipage}{0.875\linewidth}
En Afrique subsaharienne, la transformation digitale en santé se déploie dans un contexte singulier. Les systèmes de santé y  sont caractérisés par une insuffisance chronique de ressources financières, humaines et matérielles (Nabyonga- Orem\&Tashobya,2016).Les infrastructures numériques, bien que croissantes avec l'essor du téléphone mobile et d'Internet, restent inégales et parfois précaires .Les coupures d'électricité fréquentes, l'instabilité des réseaux et la faible capacité technique du personnel constituent des contraintes majeures.
\end{minipage}
\end{center}

